\section{Virtual Lab：多 Agent 科研组织的最小闭环}

Virtual Lab 是本讲第一部分，也是最像``真实研究组数字孪生''的系统。它并不假设一个全能 Agent，而是把科研活动拆成角色协作网络：PI/Professor Agent 负责任务分解与团队配置，多个 Student Agent 按学科专长并行推进，Critic Agent 作为系统内审查者持续挑错。

图\ref{fig:vl-design} 展示了 Virtual Lab 的三层结构：Agent 创建、团队会议、一对一会议。\footnote{来源：Berkeley RDI 课程视频 \url{https://www.youtube.com/watch?v=yqPIsTTdUkc}，关键帧时间戳 00:05:32。}

\begin{figure}[H]
\centering
\includegraphics[width=0.95\textwidth]{frame-06-pi-agent-flow.jpg}
\caption{Virtual Lab 结构图：Agent 创建、Team Meeting、Individual Meeting}
\label{fig:vl-design}
\end{figure}

\begin{knowledgebox}{读图：三层会议结构承担什么计算}
最上层根据课题创建互补角色；中层把多条假设放进同一组会做交叉质询；下层让单个专家与 PI 深挖高潜力分支。图中真正重要的不是 Agent 数量，而是信息何时广播、何时局部化、何时由 PI 收敛。若所有讨论都进入同一长上下文，角色分工会退化成昂贵的重复生成。
\end{knowledgebox}

\subsection{角色定义与任务路由}

在 Zou 的描述里，人类研究者通常先与 PI Agent 对话，PI 再决定需要哪些子专家。这意味着系统有一个\textbf{任务路由层}：同样是``生物医药''，不同课题可能需要不同组合，例如 immunology + structure modeling + optimization，或者 genomics + statistics + causal inference。

这一层的意义在于，复杂科研任务并不是一次 prompt 能完成的线性过程，而是带有``先验不完整''和``目标随中间结果更新''特征的动态流程。PI Agent 不是只做拆分，还做\textbf{团队重构}：随着新证据出现，替换或追加 specialist。

\begin{knowledgebox}{PI Agent 的工程职责可以抽象为三件事}
\begin{itemize}
    \item \textbf{Capacity Planning}：为当前课题配置足够、但不过度的专家角色。
    \item \textbf{Interface Contract}：定义每个 Agent 输入输出格式，减少讨论漂移。
    \item \textbf{Iteration Control}：决定何时继续探索、何时收敛、何时升级给人类。
\end{itemize}
\end{knowledgebox}

\subsection{会议机制：并行探索与集中汇总}

Virtual Lab 的会议不是社交形式，而是计算编排策略。组会让多个 Agent 并行提出假设，单聊用于把某条分支做深，再回到组会做 cross-check。这个机制本质上是在\textbf{宽搜（breadth-first）与深搜（depth-first）之间切换}。

Zou 讲到``AI 团队会议可以在几秒到几分钟内完成''，这使得许多传统团队难以承担的对照实验（不同参数设定、不同先验假设、不同人格配置）可以低成本并行执行。换言之，Virtual Lab 不是``自动写结论''，而是``自动做研究过程中的大量中间尝试''。

\begin{importantbox}{会议编排带来的三个直接收益}
\begin{itemize}
    \item \textbf{探索覆盖率提高}：并行分支数量显著增加。
    \item \textbf{沟通损耗降低}：机器间协议化交互替代口语化反复解释。
    \item \textbf{可审计性增强}：每轮讨论、每次工具调用都可追踪。
\end{itemize}
\end{importantbox}

\subsection{系统行为中的``社会动力学''}

课程中特别提到``Agents have their own social dynamics''。这不是噱头，指的是当多个 Agent 拥有不同偏好、不同不确定性处理方式时，系统会出现类似真实团队的行为：有人保守、有人冒进、有人倾向先做可行解、有人倾向追求高潜力路线。

如果没有治理机制，这些``社会动力学''会导致议题漂移、局部共识锁死或重复讨论。Virtual Lab 通过 Critic、会议摘要、PI 汇总三层机制，把动力学从噪声转化为可用信号。

\begin{warningbox}{仅靠多数投票不等于高质量科研决策}
科研里的``正确''常常是稀疏信号，早期阶段可能只被少数分支捕捉。若系统只做多数投票，会系统性压制少数高价值路径。Virtual Lab 的关键是保留异议并结构化记录，而不是快速投票定稿。
\end{warningbox}

\subsection{本章小结}

Virtual Lab 的核心价值在组织层：PI 路由、并行会议、单聊深挖、Critic 纠偏、最终汇总。它把科研活动从``个人脑内流程''转成``可编排、可追踪、可复盘''的系统流程。

